\section{Conclusion}

Starting from THUIR@COLIEE 2023, this paper re-examines long-context handling, query representation, negative sampling, and downstream ranking integration in legal case retrieval. We find that when the whole judgment is used as the query, ModernBERT + DAP + SFT reaches 0.1650 F1@5 on COLIEE 2026 validation, outperforming BM25 at 0.1499 and SAILER-cos at 0.1213. In contrast, when the query is replaced by reference sentence extraction, both neural models degrade substantially. This shows that the value of long-context encoders depends on actually preserving enough legal context in the query.

We also show that long-context encoders are not plug-and-play solutions. ModernBERT without SFT is almost unusable for legal case retrieval, and static BM25 hard negatives make training nearly collapse. The system only starts to train reliably after introducing dynamic negative sampling, year-based scope filtering, and legal-domain continued pretraining.

Finally, dense retrieval is not the end of the pipeline. Once lexical, dense, length, structural, year, and chunk signals are fused with LightGBM learning to rank and followed by post-processing, the final FLNLP submission reaches 0.2935 F1 on the official COLIEE 2026 Task~1 leaderboard and ranks 7th among 22 teams. Overall, the main challenge of legal dense retrieval is not simply choosing a larger model; it is enabling the model to see long documents under realistic hardware constraints, keeping the right query representation, and continuously providing informative training and ranking signals.
